\documentclass[10pt,a4paper]{amsart}
\usepackage{amsmath,amssymb,amsthm,graphicx}
\usepackage{verbatim}
\newenvironment{reference}{\comment}{\endcomment}
%\newenvironment{reference}{}{}
\newenvironment{slogan}{\comment}{\endcomment}
\newenvironment{history}{\comment}{\endcomment}

% For commutative diagrams we use Xy-pic
\usepackage[all]{xy}

% We use 2cell for 2-commutative diagrams.
\xyoption{2cell}
% Theorem environments.
%
\theoremstyle{plain}
\newtheorem{theorem}[subsection]{Theorem}
\newtheorem{proposition}[subsection]{Proposition}
\newtheorem{lemma}[subsection]{Lemma}

\theoremstyle{definition}
\newtheorem{definition}[subsection]{Definition}
\newtheorem{example}[subsection]{Example}
\newtheorem{exercise}[subsection]{Exercise}
\newtheorem{situation}[subsection]{Situation}

\theoremstyle{remark}
\newtheorem{remark}[subsection]{Remark}
\newtheorem{remarks}[subsection]{Remarks}

\numberwithin{equation}{subsection}

% Macros
%
\def\lim{\mathop{\mathrm{lim}}\nolimits}
\def\colim{\mathop{\mathrm{colim}}\nolimits}
\def\Spec{\mathop{\mathrm{Spec}}}
\def\Hom{\mathop{\mathrm{Hom}}\nolimits}
\def\Ext{\mathop{\mathrm{Ext}}\nolimits}
\def\SheafHom{\mathop{\mathcal{H}\!\mathit{om}}\nolimits}
\def\SheafExt{\mathop{\mathcal{E}\!\mathit{xt}}\nolimits}
\def\Sch{\mathit{Sch}}
\def\Mor{\mathop{\mathrm{Mor}}\nolimits}
\def\Ob{\mathop{\mathrm{Ob}}\nolimits}
\def\Sh{\mathop{\mathit{Sh}}\nolimits}
\def\NL{\mathop{N\!L}\nolimits}
\def\CH{\mathop{\mathrm{CH}}\nolimits}
\def\proetale{{pro\text{-}\acute{e}tale}}
\def\etale{{\acute{e}tale}}
\def\QCoh{\mathit{QCoh}}
\def\Ker{\mathop{\mathrm{Ker}}}
\def\Im{\mathop{\mathrm{Im}}}
\def\Coker{\mathop{\mathrm{Coker}}}
\def\Coim{\mathop{\mathrm{Coim}}}

% Boxtimes
%
\DeclareMathSymbol{\boxtimes}{\mathbin}{AMSa}{"02}

%
% Macros for moduli stacks/spaces
%
\def\QCohstack{\mathcal{QC}\!\mathit{oh}}
\def\Cohstack{\mathcal{C}\!\mathit{oh}}
\def\Spacesstack{\mathcal{S}\!\mathit{paces}}
\def\Quotfunctor{\mathrm{Quot}}
\def\Hilbfunctor{\mathrm{Hilb}}
\def\Curvesstack{\mathcal{C}\!\mathit{urves}}
\def\Polarizedstack{\mathcal{P}\!\mathit{olarized}}
\def\Complexesstack{\mathcal{C}\!\mathit{omplexes}}
% \Pic is the operator that assigns to X its picard group, usage \Pic(X)
% \Picardstack_{X/B} denotes the Picard stack of X over B
% \Picardfunctor_{X/B} denotes the Picard functor of X over B
\def\Pic{\mathop{\mathrm{Pic}}\nolimits}
\def\Picardstack{\mathcal{P}\!\mathit{ic}}
\def\Picardfunctor{\mathrm{Pic}}
\def\Deformationcategory{\mathcal{D}\!\mathit{ef}}
\usepackage{cleveref}
\crefname{originallemma}{original lemma}{original lemmas}
\Crefname{originallemma}{Original lemma}{Original lemmas}
\makeatletter
\newlabel{lemma-characterize-zero-local}{{00HN}{source}{}{sourcepointer-1}{}}
\newlabel{lemma-characterize-zero-local@cref}{{[originallemma][0][]00HN}{[1][0][]source}{}{}{}}
\makeatother
\crefname{originallemma}{original lemma}{original lemmas}
\Crefname{originallemma}{Original lemma}{Original lemmas}
\makeatletter
\newlabel{lemma-universal-property-localization-module}{{07K0}{source}{}{sourcepointer-2}{}}
\newlabel{lemma-universal-property-localization-module@cref}{{[originallemma][0][]07K0}{[1][0][]source}{}{}{}}
\makeatother
\crefname{originallemma}{original lemma}{original lemmas}
\Crefname{originallemma}{Original lemma}{Original lemmas}
\makeatletter
\newlabel{proposition-localize-twice-module}{{02C7}{source}{}{sourcepointer-3}{}}
\newlabel{proposition-localize-twice-module@cref}{{[originallemma][0][]02C7}{[1][0][]source}{}{}{}}
\makeatother
\title{The two-term Cech sequence of a module for a covering by basic opens is exact}
\author{The Stacks Project Authors}
\date{stacks-project:algebra.tex:local-snapshot-20261012; GNU FDL 1.2 or later}
\begin{document}
\maketitle
\noindent\textit{Editorial note.} All delivered mathematical text is the author's own, copied from the pinned Stacks Project chapter file algebra.tex; no statement or proof step was rewritten, reordered or removed.\par\medskip
\noindent\textit{Editorial note.} Preamble provenance, as required for this source: the frozen snapshot stages the chapter's own preamble file beside it (preamble.tex, 217 lines, 8116 bytes, the same shared Stacks-wide preamble the published chow and etale-cohomology cards use); the first line of the pinned chapter file is \textbackslash{}input\{preamble\}, so no preamble text lives inside the chapter itself. The card takes that file's lines 8-18 (the verbatim package and the three \textbackslash{}comment-based environments reference/slogan/history at lines 9-12, section 13.2, together with the Xy-pic setup) and its lines 154 to the runtime end of the file (the native theorem declarations and mathematical macros). Not one character was changed, no package was added, no macro was defined, and the excerpt serves only as preamble, never as a span coordinate.\par\medskip
\noindent\textit{Editorial note.} The same-chapter pointers are located at their \textbackslash{}label lines inside the pinned chapter file; the cross-chapter pointers are located in the project's own tag list staged in the snapshot (sources/stacks-chow/tags.txt), at the verbatim lines carrying their canonical tags.\par\medskip
\noindent\textit{Editorial note.} The counted proof refers to a same-chapter exactness criterion by label; it is recorded as a navigation stub pointed at its \textbackslash{}label line, and the argument consuming it stays inside the counted span.\par\medskip
\noindent\textit{Editorial note.} The counted spans contain no \textbackslash{}cite, so the card carries no bibliography.\par\medskip
\noindent\textit{Allowed standard prerequisites (editorial).}\begin{itemize}
\item Exactness can be checked after localising at maximal ideals. Invoked at the start of the counted proof (lines 4540-4542). Reduces the exactness of the displayed sequence to the local case.
\item Localising a localisation by a product of the covering elements. Used throughout the local computation (lines 4543-4567). Simplifies the terms of the sequence once one of the covering elements is invertible.
\end{itemize}\medskip
\ifdefined\hypertarget\hypertarget{sourcepointer-1}{}\fi
\noindent\textit{Source reference.} Original lemma 00HN: sources/stacks-algebra/pinned-original-source.tex, source lines 4268--4268. Navigation pointer for \textbackslash{}ref\{lemma-characterize-zero-local\} used inside the counted span (line 4542); the argument that consumes it is retained in the span.\par\medskip
\ifdefined\hypertarget\hypertarget{sourcepointer-2}{}\fi
\noindent\textit{Source reference.} Original lemma 07K0: sources/stacks-algebra/pinned-original-source.tex, source lines 1388--1388. Navigation pointer for \textbackslash{}ref\{lemma-universal-property-localization-module\} used inside the counted span (line 4554); the argument that consumes it is retained in the span.\par\medskip
\ifdefined\hypertarget\hypertarget{sourcepointer-3}{}\fi
\noindent\textit{Source reference.} Original lemma 02C7: sources/stacks-algebra/pinned-original-source.tex, source lines 1519--1519. Navigation pointer for \textbackslash{}ref\{proposition-localize-twice-module\} used inside the counted span (line 4548); the argument that consumes it is retained in the span.\par\medskip
\begin{lemma}
\label{lemma-cover-module}
Let $R$ be a ring. Let $f_1, \ldots, f_n$ be elements of $R$
generating the unit ideal. Let $M$ be an $R$-module.
The sequence
$$
0 \to
M \xrightarrow{\alpha}
\bigoplus\nolimits_{i = 1}^n M_{f_i} \xrightarrow{\beta}
\bigoplus\nolimits_{i, j = 1}^n M_{f_i f_j}
$$
is exact, where $\alpha(m) = (m/1, \ldots, m/1)$
and $\beta(m_1/f_1^{e_1}, \ldots, m_n/f_n^{e_n})
= (m_i/f_i^{e_i} - m_j/f_j^{e_j})_{(i, j)}$.
\end{lemma}

\begin{proof}
It suffices to show that the localization of the sequence at
any maximal ideal $\mathfrak m$ is exact, see
Lemma \ref{lemma-characterize-zero-local}.
Since $f_1, \ldots, f_n$ generate the unit ideal,
there is an $i$ such that $f_i \not \in \mathfrak m$.
After renumbering we may assume $i = 1$.
Note that $(M_{f_i})_\mathfrak m = (M_\mathfrak m)_{f_i}$
and $(M_{f_if_j})_\mathfrak m = (M_\mathfrak m)_{f_if_j}$, see
Proposition \ref{proposition-localize-twice-module}.
In particular $(M_{f_1})_\mathfrak m = M_\mathfrak m$ and
$(M_{f_1 f_i})_\mathfrak m = (M_\mathfrak m)_{f_i}$, because
$f_1$ is a unit.
Note that the maps in the sequence are the canonical ones
coming from
Lemma \ref{lemma-universal-property-localization-module}
and the identity map on $M$.
Having said all of this, after replacing $R$ by $R_\mathfrak m$,
$M$ by $M_\mathfrak m$, and $f_i$ by their image in $R_\mathfrak m$,
and $f_1$ by $1 \in R_\mathfrak m$,
we reduce to the case where $f_1 = 1$.

\medskip\noindent
Assume $f_1 = 1$. Injectivity of $\alpha$ is now trivial. Let
$m = (m_i) \in \bigoplus_{i = 1}^n M_{f_i}$ be in the kernel of $\beta$.
Then $m_1 \in M_{f_1} = M$. Moreover, $\beta(m) = 0$
implies that $m_1$ and $m_i$ map to the same element of
$M_{f_1f_i} = M_{f_i}$. Thus $\alpha(m_1) = m$ and the
proof is complete.
\end{proof}

\end{document}
